\section{The full small Dolbeault spectrum}\label{sec:spectrum}

We now compare the elliptic operator with the finite form. The first
step holds for every sufficiently small simple coefficient array,
whether balanced or not. All Hilbert spaces in this section are
complex, and all endomorphisms are pointwise trace-free.

\begin{proposition}\label{prop:schur}
Let $Z'$ be small with scalar centralizer. Use the split metric
$h_{0,s}$ on $E_{s,Z'}$, and let $\nu_j(s,Z')$ be the eigenvalues
of \eqref{eq:finite-form}. There are constants $C,\gamma_1>0$,
uniform in the fixed background neighborhood, such that the
ordered Dolbeault eigenvalues $\widetilde\lambda_j$ satisfy
\begin{equation}\label{eq:schurbound}
(1-Cq_s(Z'))\nu_j(s,Z')\le\widetilde\lambda_j
\le\nu_j(s,Z')\quad(1\le j\le N),\qquad
\widetilde\lambda_{N+1}\ge\gamma_1/2.
\end{equation}
In particular exactly $N$ eigenvalues lie below a fixed positive
threshold for sufficiently small $Z'$.
\end{proposition}

\begin{proof}
Write $D_0=\dbar_{\End G}$ and
$D=D_0+\mathcal A$, where $\mathcal A=[N_s(Z'),\,\cdot\,]$.
The kernel of $D_0$ is the full parallel space $\HH$ of
\eqref{eq:finite-space}, by the fixed cross-Hom vanishing.
Let $P$ be complex $L^2$ projection onto $\HH$, and put $Q=1-P$.
These are the complex spectral projections, distinct from the
real metric-coordinate projections of Section~\ref{sec:curvature}.
The positive spectrum of $D_0^*D_0$ on $Q$ has a uniform lower
bound $\gamma_0>0$.

For every parallel $X\in\HH$, $\mathcal AX$ is harmonic,
because it is a constant multiplicity combination of the harmonic
Hom forms. Thus
\begin{equation}\label{eq:harmoniccancel}
D_0^*\mathcal AX=0,\qquad
V_s^{-1}\|DX\|_{L^2}^2=\mathcal Q_{s,Z'}(X),\qquad
V_s^{-1}\|X\|_{L^2}^2=\|X\|_r^2.
\end{equation}
Let $q=q_s(Z')$. Since the harmonic coefficient space is fixed and
finite-dimensional, its $C^0$ and $L^2$ norms are uniformly
equivalent. Consequently
$\|\mathcal A\|_{L^2\to L^2(0,1)}\le C\sqrt q$.
For $y\in QH^1$ this gives
\[
\|Dy\|_{L^2}\ge(\sqrt{\gamma_0}-C\sqrt q)\|y\|_{L^2}.
\]
The complementary compression $C_0=QD^*DQ$, understood initially
as a closed quadratic form, therefore satisfies $C_0\ge\gamma_1I$
for some fixed $\gamma_1>0$ and all sufficiently small $q$.

Set $G_0=PD^*DP$ and $B=QD^*DP$. Harmonicity in
\eqref{eq:harmoniccancel} gives, as equalities of smooth sections,
\begin{equation}\label{eq:relativecoupling}
BX=Q\mathcal A^*\mathcal AX,\qquad
\|BX\|_{L^2}^2\le Cq\|\mathcal AX\|_{L^2}^2
=Cq\langle X,G_0X\rangle_{L^2}.
\end{equation}
The right side is each vector's own finite energy. There is no
replacement by $q\|X\|^2$ and no division by $\nu_1$.

We specify the form domains to justify the spectral comparison.
The form domain of $D^*D$ is $H^1(\End_0G)$ and its operator
domain is $H^2(\End_0G)$. Removing the finite-dimensional smooth
space $\HH$ preserves these Sobolev spaces; the form and operator
domains of $C_0$ are $QH^1$ and $QH^2$. In particular $B$ maps
the finite space into $L^2$, and $(C_0-\lambda)^{-1}$ maps $QL^2$
into $QH^2$ for $0\le\lambda\le\gamma_1/2$.

For such $\lambda$ define the finite Schur complement
\begin{equation}\label{eq:schur}
S_\lambda=G_0-\lambda I-B^*(C_0-\lambda)^{-1}B.
\end{equation}
Equation~\eqref{eq:relativecoupling} gives the Hermitian-form bounds
\begin{equation}\label{eq:schurforms}
(1-Cq)G_0-\lambda I\le S_\lambda\le G_0-\lambda I.
\end{equation}
Completion of the square on $\HH\oplus QH^1$ gives
\begin{align}
\mathfrak q_{D^*D-\lambda}(X+y)
={}&\langle S_\lambda X,X\rangle\notag\\
&+\left\|(C_0-\lambda)^{1/2}
\bigl(y+(C_0-\lambda)^{-1}BX\bigr)\right\|^2.\label{eq:inertia}
\end{align}
The triangular change of variables here is a bijection of the
form domain, since its correction lies in $QH^2$. The negative
index of $D^*D-\lambda$ thus equals the negative index of
$S_\lambda$.

The norm of $G_0$ is $O(q)$. At $\lambda=\gamma_1/2$ the
finite Schur complement is strictly negative, whereas the
complementary operator is positive. Exactly $N$ eigenvalues
therefore lie below this threshold and the threshold is not an
eigenvalue. For the indexed lower bound put $\alpha=1-Cq>0$.
If $\lambda<\alpha\nu_j$, the form $\alpha G_0-\lambda I$
has at most $j-1$ negative directions. The same is true of
$S_\lambda$ by \eqref{eq:schurforms}; \eqref{eq:inertia}
then implies $\widetilde\lambda_j\ge\alpha\nu_j$ by passage
to the limit. For the upper bound apply ordinary min--max to
the first $j$ finite eigenvectors inside $\HH$.
These arguments also apply when eigenvalues coincide. They
establish \eqref{eq:schurbound} with complex multiplicities.
\end{proof}

\begin{lemma}\label{lem:gaugespectrum}
Let $g=e^{Y/2}$ for $Y\in\hh$, put $a=\|Y\|_{\op}$, and let
$Z'=g\cdot Z$. For all $1\le j\le N$,
\begin{equation}\label{eq:gaugecomparison}
e^{-4a}\nu_j(s,Z)\le\nu_j(s,Z')\le e^{4a}\nu_j(s,Z).
\end{equation}
This comparison holds for arbitrary positive harmonic Gram matrices.
\end{lemma}

\begin{proof}
Conjugation $T_gX=(g_vX_vg_v^{-1})_v$ preserves the full complex
space $\HH$, since it preserves each block trace. The bounds
$\|g_v^{\pm1}\|_{\op}\le e^{a/2}$ imply
\begin{equation}\label{eq:gaugemass}
e^{-2a}\|X\|_r^2\le\|T_gX\|_r^2\le e^{2a}\|X\|_r^2.
\end{equation}
The exact intertwining identity
\[
[T_gX,g\cdot Z]_{ij,a}=g_i[X,Z]_{ij,a}g_j^{-1}
\]
gives the same factors for the squared energy. For a nondiagonal
Gram matrix, apply its positive square root to the arrow index.
Left and right multiplicity operations commute with this change
because every arrow in the block has the same source and target.
Thus the same energy bounds hold without changing constants.
Rayleigh quotients have factors $e^{\pm4a}$, and generalized
min--max through the invertible complex map $T_g$ proves the lemma.
\end{proof}

The map $T_g$ need not preserve Hermitian matrices. Its use on
$\HH$ is one reason to retain the full complex operator throughout
the argument. At exact central balance there is an additional
adjoint symmetry, discussed in Section~\ref{sec:scope}, but no
such symmetry is needed during this comparison.

\begin{proof}[Proof of Theorem~\ref{thm:main}]
Proposition~\ref{prop:actual-stability} proves the stability
equivalence and metric statement. At a stable parameter choose
the independently balanced $Z$. Theorem~\ref{thm:calibration}
gives $Z'=e^{Y/2}\cdot Z$ with $\|Y\|_{\op}\le A\sigma$
and $\|u\|_{C^0}\le C\sigma$. Moreover
$q_s(Z')\le e^{2A\sigma}\sigma$.
Proposition~\ref{prop:schur} and Lemma~\ref{lem:gaugespectrum}
therefore compare the split-metric small eigenvalues at $Z'$
to $\nu_j(s,Z)$ with a common factor $1+O(\sigma)$.

It remains to pass to the actual metric on the same holomorphic
bundle $E_{s,Z'}$. If two bundle metrics satisfy
$e^{-c}h_0\le h_1\le e^ch_0$, their induced squared norms on
$\End G$ lie between $e^{-2c}$ and $e^{2c}$ times each other.
The same is true on $\End G$-valued $(0,1)$-forms, since the
base metric is fixed in this comparison. The differential $D$
is unchanged. Hence every Rayleigh quotient has comparison
factors $e^{\pm4c}$ on the common $H^1$ form domain. Apply this
with $h_1=h_{0,s}e^u$ and $c=C\sigma$.

The trace-free spaces in this argument are invariant for both
Laplacians. Indeed $\End G=\C I\oplus\End_0G$ is orthogonal
for every induced endomorphism metric, $D(fI)=(\dbar f)I$, and
$\tr(D\xi)=\dbar\tr\xi$. This orthogonal splitting also
holds on one-forms. Both $D$ and its adjoint preserve the
corresponding summands. Min--max therefore applies to the same
trace-free domain for each metric.

The additional factor is $1+O(\sigma)$, so it proves
\eqref{eq:main-spectrum} for $1\le j\le N$. It also leaves a
fixed positive fraction of the complementary bound in
Proposition~\ref{prop:schur}. Finally holomorphic pullback,
with the metric pulled back at the same time, is unitary on
the relevant $L^2$ spaces and intertwines the Dolbeault
operators. It identifies these eigenvalues with those on
the original $E$, regardless of the size of the original
balancing gauge. Shrinking $\mathcal U$ absorbs all uniform
exponential factors into the stated constants.
\end{proof}
